% GENERATED FILE. Edit canonical lessons, then run npm run generate:books.
% canonical-source-hash: fnv1a64:181803044a7f3764
% canonical-lessons: ES-C368-abeja, ES-C368-mosca, ES-C368-cerdo, ES-C368-lobo, ES-C368-toro

\chapter{Five Animals, Two Of Them Small}
\label{ch:la-abeja-y-la-mosca}
\hlchaptermodality{368}

\begin{chapteropening}
\textbf{By the end of this chapter:} \emph{I can say la abeja, la mosca, el cerdo, el lobo and el toro, and name five animals a field might hold.}

Complete ES-C368-toro: name five animals, two of them small enough to sit on a hand.
\end{chapteropening}

\section[la abeja]{la abeja — the little one that outlived the big one}

\label{lesson:ES-C368-abeja}



\begin{quote}
Say \emph{the potato}, then \emph{sugar}. (\emph{La patata}; \emph{el azúcar}.)
\end{quote}

\subsection*{What to know first}
\begin{itemize}
\raggedright
  \item la miel — what she makes.
  \item la flor — where she works.
\end{itemize}

\begin{sounds}
\begin{itemize}
\raggedright
  \item \emph{a-BE-ja}, three beats with the weight on the second. The \emph{b} between vowels is soft, the lips barely meeting, and the \emph{j} scrapes at the back. $\to$ pronunciation reference
\end{itemize}
\end{sounds}

\begin{cousinweb}
\textbf{la abeja} — \textquotedblleft{}the bee.\textquotedblright{}

Latin had a perfectly good word for a bee: \textbf{apis}. Spanish does not use it. \emph{Abeja} comes from \textbf{apicula}, which is \emph{apis} with a diminutive ending — \textquotedblleft{}the little bee.\textquotedblright{}

That happened more than once, and once you notice it you will keep meeting it:

\begin{itemize}
\raggedright
  \item \textbf{oveja}, a sheep, comes from \emph{ovicula}, the little \emph{ovis};
  \item \textbf{oreja}, an ear, comes from \emph{auricula}, the little \emph{auris}.
\end{itemize}

In each pair the plain Latin noun died in Spanish and the affectionate small form took over the ordinary job. Whatever the Romans were doing when they talked about their animals and their bodies, they were doing it in diminutives, and the diminutive is what survived.

One word that looks like it belongs in that list does not. \textbf{Conejo}, a rabbit, comes from Latin \emph{cunīculus}, which \textbf{has} the diminutive ending but no base noun anywhere behind it — Latin most likely borrowed the whole word from an older language of Spain. It is a shape, not a formation.

English went the other way and took the plain \textbf{apis} for its bookish words: an \textbf{apiary} is a bee-yard, \textbf{apiculture} is bee-keeping. So one Latin noun ended up split between two languages, Spanish holding the small form and English the large one.

English \textbf{bee} is a different word entirely, Germanic, with no link to \emph{apis}. And \emph{apis} itself is a puzzle: it has no convincing ancestry in the older layers of Latin's own family, and looks like a word the Romans picked up from someone who was in Italy before them.
\end{cousinweb}

\begin{grammarlens}[title={What you've built}]
You can now name the maker as well as what she makes: \emph{la abeja hace la miel}.
\end{grammarlens}

\subsection*{Your turn}
Say these aloud:

\begin{itemize}
\raggedright
  \item \emph{la abeja}
  \item \emph{una abeja}
  \item \emph{la abeja hace la miel}
\end{itemize}

\begin{tcolorbox}[breakable,colback=teal!4,colframe=teal!35!black,title={Before you move on}]
What does \emph{abeja} come from, \emph{apis} or \emph{apicula}? (\emph{Apicula}, the little one.) And what is an apiary? (A place for bees.)
\end{tcolorbox}
\section[la mosca]{la mosca — the fly that became a gun}

\label{lesson:ES-C368-mosca}



\begin{quote}
Say \emph{sugar}, then \emph{the bee}. (\emph{El azúcar}; \emph{la abeja}.)
\end{quote}

\subsection*{What to know first}
\begin{itemize}
\raggedright
  \item la comida — what a \emph{mosca} comes for.
  \item la sopa — the classic place to find one.
\end{itemize}

\begin{sounds}
\begin{itemize}
\raggedright
  \item \emph{MOS-ca}, two beats with the weight on the first. The \emph{sc} is a clean \emph{s} against a hard \emph{k}, with no vowel between them. $\to$ pronunciation reference
\end{itemize}
\end{sounds}

\begin{cousinweb}
\textbf{la mosca} — \textquotedblleft{}the fly.\textquotedblright{}

Latin \textbf{musca}, straight down, no complications. The interesting part is what grew out of it.

Spanish makes small things with \textbf{-ito}, so a little \emph{mosca} is a \textbf{mosquito}. English borrowed that word entire — Spanish ending and all — which is why an English speaker is carrying a Spanish diminutive around without knowing it. Say \emph{mosquito} and you have said \textquotedblleft{}little fly\textquotedblright{} in Spanish.

Now the strange one. French made its own small form of \emph{musca} and gave it to a bird: the \textbf{sparrowhawk}, small, and speckled when it moves. Firearms of that century were routinely named after hunting birds — the \textbf{falcon} and the \textbf{falconet} were both cannon — and the little-hawk name went to a shoulder gun. English took it from French as \textbf{musket}.

So the chain runs: fly, little fly, small hawk, gun. A musket is a mosquito with a longer barrel.

One more relative, offered with a caution. Greek \emph{muia} and English \textbf{midge} are usually placed alongside \emph{musca} far back in the family tree. But a word for a small buzzing insect is exactly the kind of word a language invents afresh out of the noise the thing makes, so a resemblance here proves less than the same resemblance would prove elsewhere.
\end{cousinweb}

\begin{grammarlens}[title={What you've built}]
You can now name the household nuisance: \emph{la mosca}.
\end{grammarlens}

\subsection*{Your turn}
Say these aloud:

\begin{itemize}
\raggedright
  \item \emph{la mosca}
  \item \emph{una mosca}
  \item \emph{hay una mosca en la sopa}
\end{itemize}

\begin{tcolorbox}[breakable,colback=teal!4,colframe=teal!35!black,title={Before you move on}]
What does \emph{mosquito} mean taken apart? (Little fly.) And what animal stands between the fly and the gun? (The sparrowhawk.)
\end{tcolorbox}
\section[el cerdo]{el cerdo — the animal named after its hair}

\label{lesson:ES-C368-cerdo}



\begin{quote}
Say \emph{the bee}, then \emph{the fly}. (\emph{La abeja}; \emph{la mosca}.)
\end{quote}

\subsection*{What to know first}
\begin{itemize}
\raggedright
  \item la granja — where you meet one.
  \item la carne — what it turns into.
\end{itemize}

\begin{sounds}
\begin{itemize}
\raggedright
  \item \emph{CER-do}, two beats with the weight on the first. The \emph{c} before \emph{e} is soft, and the \emph{d} after \emph{r} is gentle, closer to the \emph{th} of English \emph{this}. $\to$ pronunciation reference
\end{itemize}
\end{sounds}

\begin{cousinweb}
\textbf{el cerdo} — \textquotedblleft{}the pig.\textquotedblright{}

This animal has two names in Spanish, and the ordinary one is the odd one.

\textbf{Cerda} is a coarse bristle — the stiff hair a pig has instead of fur. Farmers spoke of \emph{ganado de la cerda}, the bristle livestock, and over a few centuries the herd's nickname squeezed out the animal's name. A \emph{cerdo} is \textquotedblleft{}the bristly one.\textquotedblright{} Where \emph{cerda} itself came from is argued over: one account takes it to a Latin word for a tuft of hair, another to a Latin word for a stiff bristle, and neither has won.

The name it pushed aside is still standing. Latin \textbf{porcus} came down into Spanish perfectly regularly as \textbf{puerco}, which nobody has to be taught and everybody understands. \emph{Cerdo} did not replace it; it took the neutral, unremarkable slot and left \emph{puerco} the blunter one.

English holds that same \emph{porcus} twice over, arriving by two different doors:

\begin{itemize}
\raggedright
  \item \textbf{pork} came in from French, which is why English eats \emph{pork} and keeps a \emph{pig};
  \item \textbf{farrow}, a litter of piglets, was never borrowed at all. It came down the Germanic line from the same ancient word \emph{porcus} came down the Latin line.
\end{itemize}

Two languages, one old word for a pig, and in each of them a second word elbowing in beside it.
\end{cousinweb}

\begin{grammarlens}[title={What you've built}]
You can now name the farm animal and its meat: \emph{el cerdo}, \emph{la carne de cerdo}.
\end{grammarlens}

\subsection*{Your turn}
Say these aloud:

\begin{itemize}
\raggedright
  \item \emph{el cerdo}
  \item \emph{un cerdo}
  \item \emph{la carne de cerdo}
\end{itemize}

\begin{tcolorbox}[breakable,colback=teal!4,colframe=teal!35!black,title={Before you move on}]
What does \emph{cerda} mean? (A bristle.) And what is the Spanish word that comes straight from Latin \emph{porcus}? (\emph{Puerco}.)
\end{tcolorbox}
\section[el lobo]{el lobo — the wolf that came in from the next valley}

\label{lesson:ES-C368-lobo}



\begin{quote}
Say \emph{the fly}, then \emph{the pig}. (\emph{La mosca}; \emph{el cerdo}.)
\end{quote}

\subsection*{What to know first}
\begin{itemize}
\raggedright
  \item el bosque — where he lives.
  \item la oveja — what he is counted against.
\end{itemize}

\begin{sounds}
\begin{itemize}
\raggedright
  \item \emph{LO-bo}, two beats with the weight on the first. The \emph{b} between vowels is soft; the lips come close without quite closing. $\to$ pronunciation reference
\end{itemize}
\end{sounds}

\begin{cousinweb}
\textbf{el lobo} — \textquotedblleft{}the wolf.\textquotedblright{}

\emph{Lobo} is Latin \textbf{lupus}, and \emph{lupus} is where the story gets strange, because \emph{lupus} is not the shape Latin should have ended up with.

Languages change sound by sound, and the changes are regular enough that a scholar can predict what a word ought to look like. Run the ancient word for a wolf through Latin's own machinery and you do not get \emph{lupus}. You get something else. English \textbf{wolf} is the shape that comes out when the same ancient word goes through the Germanic machinery instead, and it behaves.

The usual explanation is that Rome did not inherit this word. It \textbf{borrowed} it from a neighbouring language of Italy — a related one, spoken up the road, in which that consonant turning into a \emph{p} was the ordinary outcome. Latin took its wolf from the next valley and kept it.

That is the mainstream account, not a proven one; another proposal starts the word from an old name for the marten instead and has trouble with the meaning. What everyone agrees on is the shape being wrong for a word Latin grew itself.

English also holds \emph{lupus} directly, borrowed centuries later off the page: \textbf{lupine} for anything wolf-like, and \textbf{lupus} for the illness, named for the way the rash was thought to eat at a face.
\end{cousinweb}

\begin{grammarlens}[title={What you've built}]
You can now put the two animals in one sentence: \emph{el lobo y la oveja}.
\end{grammarlens}

\subsection*{Your turn}
Say these aloud:

\begin{itemize}
\raggedright
  \item \emph{el lobo}
  \item \emph{un lobo}
  \item \emph{el lobo está en el bosque}
\end{itemize}

\begin{tcolorbox}[breakable,colback=teal!4,colframe=teal!35!black,title={Before you move on}]
Why is \emph{lupus} thought to be a borrowing? (Its shape is wrong for a word Latin inherited.) And which English word is the cousin that stayed home? (\emph{Wolf}.)
\end{tcolorbox}
\section[el toro]{el toro — the word that predates the languages that carry it}

\label{lesson:ES-C368-toro}



\begin{quote}
Say \emph{the pig}, then \emph{the wolf}. (\emph{El cerdo}; \emph{el lobo}.)
\end{quote}

\subsection*{What to know first}
\begin{itemize}
\raggedright
  \item la vaca — his counterpart.
  \item el campo — where the two of them stand.
\end{itemize}

\begin{sounds}
\begin{itemize}
\raggedright
  \item \emph{TO-ro}, two beats with the weight on the first. The \emph{t} touches the teeth and the \emph{r} is a single tap, not a roll. $\to$ pronunciation reference
\end{itemize}
\end{sounds}

\begin{cousinweb}
\textbf{el toro} — \textquotedblleft{}the bull.\textquotedblright{}

\emph{Toro} is Latin \textbf{taurus}, and \emph{taurus} is a traveller.

Something very close to it turns up in Greek, in the Celtic languages, and further east as well. That looks at first like a family resemblance — cousins who inherited the word from one ancestor. But the shapes do not line up the way inherited words line up. They agree too well in some places and not well enough in others.

The reading most scholars now take is that this is a \textbf{wandering word}: not inherited by any of them, but picked up from a language spoken around the Mediterranean before any of these languages arrived, and passed hand to hand along with the animal. A bull is a valuable thing to trade, and a traded thing brings its name.

English carries the Latin form in two ways. \textbf{Taurus} is the constellation and the zodiac sign, borrowed off the page. \textbf{Toreador} came from Spanish with the bullring attached.

And English \textbf{steer}, the ox, is sometimes filed with \emph{taurus} and sometimes kept apart, depending on whose reconstruction you follow. That disagreement is not a gap in the scholarship. It is what a word looks like when it is older than the families trying to claim it.
\end{cousinweb}

\begin{grammarlens}[title={What you've built}]
You can now name both animals in the field: \emph{el toro y la vaca}.
\end{grammarlens}

\subsection*{Your turn}
Say these aloud:

\begin{itemize}
\raggedright
  \item \emph{el toro}
  \item \emph{un toro}
  \item \emph{el toro está en el campo}
\end{itemize}

\begin{tcolorbox}[breakable,colback=teal!4,colframe=teal!35!black,title={Before you move on}]
Why is \emph{taurus} called a wandering word? (It travelled between languages instead of being inherited by them.) And what does \emph{Taurus} name in the sky? (The constellation, the bull.)
\end{tcolorbox}
